% part: intuitionistic-logic
% chapter: introduction
% section: axiomatic-derivations

\documentclass[../../../include/open-logic-section]{subfiles}

\begin{document}

\olfileid{int}{int}{axd}

\olsection{\usetoken{P}{derivation} Aksiomatis}

!!{derivation}s aksiomatis kanggo logika proposisional intuisionistik
iku sistem !!{derivation} kang paling prasaja sacara konseptual lan
pisanan sacara historis. Cara kerjane padha karo ing logika
proposisional klasik.

\begin{defn}[!!^{derivability}]
Yen $\Gamma$ iku himpunan !!{formula}s saka $\Lang L$, mula
\emph{!!{derivation}} saka $\Gamma$ iku urutan winates $!A_1$,
\dots,~$!A_n$ kang dumadi saka !!{formula}s, kanthi saben $i \le n$
nyukupi salah siji saka kahanan iki:
\begin{enumerate}
\item $!A_i \in \Gamma$; utawa
\item $!A_i$ iku aksioma; utawa
\item $!A_i$ disimpulake saka sawenehing $!A_j$ lan $!A_k$ kanthi
  $j < i$ lan $k < i$ nganggo modus ponens, yaiku
  $!A_k \ident !A_j \lif !A_i$.
\end{enumerate}
\end{defn}

\begin{defn}[Aksioma]
Himpunan $\PAx$ saka \emph{aksioma} kanggo logika proposisional
intuisionistik dumadi saka kabeh !!{formula}s kang awujud kaya mangkene:
\begin{align}
  & (!A \land !B) \lif !A \ollabel{ax:land1}\\
  & (!A \land !B) \lif !B \ollabel{ax:land2}\\
  & !A \lif (!B \lif (!A \land !B)) \ollabel{ax:land3}\\
  & !A \lif (!A \lor !B) \ollabel{ax:lor1}\\
  & !A \lif (!B \lor !A) \ollabel{ax:lor2}\\
  & (!A \lif !C) \lif ((!B \lif !C) \lif ((!A \lor !B) \lif !C)) \ollabel{ax:lor3}\\
  & !A \lif (!B \lif !A) \ollabel{ax:lif1}\\
  & (!A \lif (!B \lif !C)) \lif ((!A \lif !B) \lif (!A \lif !C)) \ollabel{ax:lif2}\\
  & \lfalse \lif !A \ollabel{ax:lfalse1}
\end{align}
\end{defn}

\begin{defn}[!!^{derivability}]
!!^a{formula}~$!A$ iku \emph{!!{derivable}} saka $\Gamma$, ditulis
$\Gamma \Proves !A$, yen ana !!a{derivation} saka $\Gamma$ kang
pungkasan ing $!A$.
\end{defn}

\begin{defn}[Teorema]
!!^a{formula}~$!A$ iku \emph{teorema} yen ana !!a{derivation}
kanggo~$!A$ saka himpunan kosong. Kita nulis $\Proves !A$ yen $!A$
iku teorema lan $\Proves/ !A$ yen dudu teorema.
\end{defn}

\begin{prop}
  Yen $\Gamma \Proves !A$ ing sistem !!{derivation} aksiomatis kanggo
  logika intuisionistik, $\Gamma \Proves !A$ uga ing logika klasik.
  Mligine, yen $!A$ iku teorema intuisionistik, uga dadi teorema klasik.
\end{prop}

\begin{proof}
  Saben aksioma intuisionistik uga aksioma klasik, mula saben
  !!{derivation} ing logika intuisionistik uga !!a{derivation} ing
  logika klasik.
\end{proof}

\end{document}
